Despite hash pre-images can be disseminated to target peers before they receive the block that contains the corresponding hashes, a robust solution is needed with withstand network partitions and peer crashes which prevent the hash pre-images from reaching all the required peers. Furthermore, if new peers are added to the Blockchain, they can only obtain the required pre-images by pulling them from peers that posses them.

To that end, prior to committing a block to the ledger, a peer iterates through all hashed key writes in all transactions of the Block and analyzes whether it needs to pull private data for each hashes key write. A hashed key write is considered as a viable candidate for pulling pre-images if:
\begin{itemize}
	\item The transaction passed validation of endorsement policies {Are we into abstracting out EPs? If yes then how should I phrase this?}. This is needed in order to prevent a denial-of-service attack in which a malicious party creates  transactions with hashed writes for which there are no pre-image in any peer in the network. Peers would try to pull these transactions and fail, hence considerably slowing down throughput.
	\item Evaluation of the collection policy on the peer's identity returns "true" {Once we define collection policies more abstractly we need to change this as well}: this is required since peers that receive pre-image requests for a certain hashed write also evaluate the same collection policy on the requester peer's identity in order to determine if that peer is part of the collection the hashed write is associated with.
	\item Whether the peer is missing the pre-image from its transient store, since by definition this means the pre-image hasn't been disseminated to the peer before the commit of the block.
	\item {Should we mention purged private data?}
\end{itemize}

Once the peer decides it needs to pull a hashed write, it needs to select potential peers which might have the pre-image it needs, and send them requests for the digests.
Upon reception of responses containing pre-images, the peer validates that the hashes on the pre-images received matches the hashed write in the block. It follows from the collision resistance property of the cryptographic hash function that if the responding peer sent a pre-image of a hashed write, then with overwhelming probability, the pre-image is identical to all pre-images in all peers in the network, hence consistency is preserved.

When sending requests for pre-images, a requester peer doesn't know up-front whether the responder peer has the pre-image, and every such "wrong guess" results in a retry which delays the block commit and thus hinders performance.

For this reason, the peers prioritize requesting a pre-image from the peers which endorsed the transaction {This is very Fabric specific, so we might want to remove this?} from which the hashed write originated. This is possible because peers constantly maintain a membership view of alive peers by periodically disseminating heartbeats. 

